\section{Languages for Regex with Positive Lookahead}
\label{sec:lookahead-semantics}
% ============================================================

This section develops the semantic foundation for regular expressions
with positive lookahead. Unlike ordinary regular expressions, whose
semantics is usually given by complete-input acceptance, positive
lookahead is evaluated at a specific input boundary. We therefore use a
window-based semantics, which records the interval of the input string on
which an expression matches.

\subsection{Regex with Positive Lookahead}

The set of regular expressions with positive lookahead over $\Sigma$, denoted by $\mathrm{REwPLA}$, is generated by the grammar in Equation~\ref{equation: grammar for standard regex}, extended with the construct $\mathsf{LA}(r)$. Unlike a standard regular expression, $\mathsf{LA}(r)$ does not consume any input string; instead, it denotes a string constraint requiring that the string following the current position belongs to the language denoted by $r$.

For example, $\mathsf{LA}((ab)^*)(a+ b)^*$ asserts at the beginning of the string that the following string has a prefix in the set $
\{\varepsilon,ab,abab,\ldots,(ab)^n,\ldots\},
$
where $n$ ranges over the natural numbers.

% ------------------------------------------------------------
\subsection{Languages with String Pairs}
\label{sec:languages-with-prefix-contexts}
% ------------------------------------------------------------


Let $\Sigma$ be an alphabet, and let $\Sigma^*$ denote the set of all finite strings over $\Sigma$. We use string pairs $(u,v)\in\Sigma^*\times\Sigma^*$, where $u$ denotes the main string and $v$ denotes the string constraint on what must follow $u$. We write $\mathcal{C}_{\Sigma}$ for the class of languages with string constraints over $\Sigma$. We define the concatenation map $\pi$ and the projection maps
$\pi_1,\pi_2$ by
$
\pi(u,v)=uv,
\:
\pi_1(u,v)=u,
\:
\pi_2(u,v)=v,
$
for every $(u,v)\in\Sigma^*\times\Sigma^*$.

We define the prefix relation $\preceq$ on strings by $
u\preceq v \:\Longleftrightarrow\: \exists z\in\Sigma^*,\ v=uz . \label{eq:prefix-relation}$ Whenever $u\preceq v$, we write $u^{-1}v$ for the unique suffix
$z\in\Sigma^*$ such that $v=uz$. Two strings are said to be prefix-compatible, written $u\bowtie v$, if one is a prefix of the other: $ u\bowtie v \:\Longleftrightarrow\: u\preceq v \ \vee\ v\preceq u . \label{eq:prefix-compatible} $

For two prefix-compatible strings $u$ and $v$ (i.e., $u \bowtie v $), we define two operators between $u$ and $v$ as follows: 
join $u\sqcup v$ as the longer one:
\begin{equation}
u\sqcup v =
\begin{cases}
v, & \text{if } u\preceq v,\\
u, & \text{if } v\preceq u,
\end{cases}
\qquad\qquad
v/u =
\begin{cases}
\varepsilon, & \text{if } v\preceq u,\\
u^{-1}v, & \text{if } u\preceq v,
\end{cases}
\label{eq:context-join and eq:context-after-consumption}
\end{equation}

where $u \sqcup v$ returns the longer string between $u$ and $v$, $v / u$ returns the left quotient of v by u if $u \preceq v$, and $\epsilon$ otherwise. 

For two prefix-constrained pairs $(u,v),(u',v')\in\Sigma^*\times\Sigma^*$, we define their partial concatenation by
\begin{equation}
(u,v)\odot(u',v')
=
\bigl(
uu',
(v/u')\sqcup v'
\bigr), \qquad v \bowtie u'v'.
\label{eq:pair-concatenation}
\end{equation}

% \begin{lemma}[\textbf{Definedness characterization}]
% \label{lem:pair-concatenation-definedness}
% The concatenation $(u,v)\odot(u',v')$ is defined if and only if there exists a string $t\in\Sigma^*$ such that $ v\preceq u't \:\text{and}\: v'\preceq t. $ and if and only if $v \bowtie u'v'$
% \end{lemma}

% \begin{proof}
% The forward direction follows immediately from
% Proposition~\ref{prop:pair-concatenation-preserves-constraints}.

% Conversely, suppose that there exists $t\in\Sigma^*$ such that
% \[
% v\preceq u't
% \qquad\text{and}\qquad
% v'\preceq t .
% \]
% Then
% \[
% u'v'\preceq u't.
% \]
% Hence both $v$ and $u'v'$ are prefixes of the same string $u't$.
% Therefore,
% \[
% v\bowtie u'v',
% \]
% and thus $(u,v)\odot(u',v')$ is defined.
% \end{proof}

\begin{proposition}
% \begin{proposition}[\textbf{Constraint preservation}]
\label{prop:pair-concatenation-preserves-constraints}
$
(u,v)\odot(u',v')=(uu',t)
\ \text{is defined}
\ \Longrightarrow\
v\preceq u't
\ \land\
v'\preceq t . \label{eq:least-residual-constraint}
$
\end{proposition}

\begin{proof}
By definition, $t=(v/u')\sqcup v'.$ We distinguish two cases. If $v\preceq u'$, then $v/u'=\varepsilon$, and hence $t=\varepsilon\sqcup v'=v'.$ Therefore, $v\preceq u'\preceq u't$ and trivially $v'\preceq t.$

Otherwise, $u'\preceq v$. Hence there exists a unique $w\in\Sigma^*$ such that $v=u'w.$ By definition, $v/u'=w.$
Since the concatenation is defined, the join
$w\sqcup v'$ is defined, and therefore $
w\preceq w\sqcup v'
\:\text{and}\:
v'\preceq w\sqcup v'.$
Since $t=w\sqcup v'$, we obtain $v=u'w\preceq u't$ and $v'\preceq t.$ Thus both constraints are preserved.
\end{proof}

\begin{lemma}
% \begin{lemma}[\textbf{Definedness characterization}]
\label{lem:pair-concatenation-definedness}
For any $(u,v),(u',v')\in\Sigma^*\times\Sigma^*$,
% \begin{equation}
$
(u,v)\odot(u',v')\text{ is defined}
\:\Longleftrightarrow\:
\exists t\in\Sigma^*.\; v\preceq u't \land v'\preceq t
$

$
\Longleftrightarrow\:
v\bowtie u'v'.
\label{eq:pair-concatenation-definedness}
$
\end{lemma}

\begin{proof}
The implication from definedness to the existence of $t$ follows
immediately from
Proposition~\ref{prop:pair-concatenation-preserves-constraints}.

Conversely, suppose that there exists $t\in\Sigma^*$ such that
$v\preceq u't$ and $v'\preceq t$. Since $v'\preceq t$, we have
$u'v'\preceq u't$. Thus, both $v$ and $u'v'$ are prefixes of the same
string $u't$, and hence $v\bowtie u'v'$. By
Definition~\eqref{eq:pair-concatenation}, the concatenation is therefore
defined.

Finally, by Definition~\eqref{eq:pair-concatenation},
$v\bowtie u'v'$ is exactly the condition under which
$(u,v)\odot(u',v')$ is defined.
\end{proof}

\begin{proposition}
% \begin{proposition}[\textbf{Least residual constraint}]
\label{prop:pair-concatenation-least-constraint}
If $(u,v)\odot(u',v')=(uu',t)$ is defined, then $t$ is the least string with respect to $\preceq$ such that $ v\preceq u't \:\text{and}\: v'\preceq t.$
\end{proposition}

\begin{proof}
By Proposition~\ref{prop:pair-concatenation-preserves-constraints}, $t$ satisfies $ v\preceq u't \: \text{and}\: v'\preceq t .$ It remains to prove that $t$ is the least such string. By definition, $t=(v/u')\sqcup v'.$ Let $s\in\Sigma^*$ be any string satisfying $v\preceq u's \:\text{and}\: v'\preceq s .$

We distinguish two cases. If $v\preceq u'$, then $v/u'=\varepsilon$, and hence $ t=\varepsilon\sqcup v'=v'. $ Since $v'\preceq s$, we immediately obtain $ t\preceq s. $ 

Otherwise, $u'\preceq v$. Write $ v=u'w, $ where $w=v/u'$. Since $ v\preceq u's, $ we have $ u'w\preceq u's. $ By left cancellation in the free monoid $\Sigma^*$, $w\preceq s.$ Together with $v'\preceq s$, both $w$ and $v'$ are prefixes of $s$. Since $ t=w\sqcup v'$ is the longer of the two prefix-compatible strings $w$ and $v'$, it follows that $ t\preceq s.$

Therefore, every string satisfying \eqref{eq:least-residual-constraint} has $t$ as a prefix. Hence $t$ is the least such string with respect to $\preceq$.
\end{proof}

\begin{remark}
Propositions~\ref{prop:pair-concatenation-preserves-constraints}
and~\ref{prop:pair-concatenation-least-constraint} characterize the
semantics of our pair concatenation. Proposition~\ref{prop:pair-concatenation-preserves-constraints}
guarantees that the resulting pair preserves the constraints imposed by both
operands. Proposition~\ref{prop:pair-concatenation-least-constraint}
further shows that $t$ is the least residual constraint, with respect
to the prefix order, satisfying these two requirements.

\end{remark}

\begin{proposition}[\textbf{Associativity}]
\label{prop:pair-concatenation-associative}
For any $(u,v),(u',v'),(u'',v'')\in\Sigma^*\times\Sigma^*$,
\begin{equation}
\bigl((u,v)\odot(u',v')\bigr)\odot(u'',v'')
=
(u,v)\odot\bigl((u',v')\odot(u'',v'')\bigr),
\label{eq:pair-concatenation-associative}
\end{equation}
whenever either side is defined. Moreover, one side is defined if and
only if the other side is defined.
\end{proposition}

\begin{proof}
Let
\[
\mathcal{T}
=
\{\,t\in\Sigma^*
\mid
v\preceq u'u''t,\;
v'\preceq u''t,\;
v''\preceq t
\,\}.
\]
We show that both parenthesizations are defined exactly when
$\mathcal{T}\neq\emptyset$, and that, whenever defined, both return
the least element of $\mathcal{T}$ as their residual constraint.

Consider first the left-associated concatenation. Let
$(u,v)\odot(u',v')=(uu',t_1)$. By
Proposition~\ref{prop:pair-concatenation-least-constraint}, $t_1$ is
the least string satisfying $v\preceq u't_1$ and $v'\preceq t_1$.
By Lemma~\ref{lem:pair-concatenation-definedness},
$(uu',t_1)\odot(u'',v'')$ is defined iff there exists a string $t$
such that $t_1\preceq u''t$ and $v''\preceq t$.

If such a $t$ exists, then
$v\preceq u't_1\preceq u'u''t$ and
$v'\preceq t_1\preceq u''t$, and hence $t\in\mathcal{T}$.
Conversely, if $t\in\mathcal{T}$, then
$v\preceq u'(u''t)$ and $v'\preceq u''t$. By the leastness of $t_1$,
we have $t_1\preceq u''t$. Together with $v''\preceq t$, this implies
that the second concatenation is defined. Hence the left-associated
concatenation is defined iff $\mathcal{T}\neq\emptyset$.

Moreover, by Proposition~\ref{prop:pair-concatenation-least-constraint},
its final residual is the least string $t$ satisfying
$t_1\preceq u''t$ and $v''\preceq t$. By the argument above, these
conditions are equivalent to $t\in\mathcal{T}$. Therefore, the final
residual is precisely the least element of $\mathcal{T}$.

Now consider the right-associated concatenation. Let
$(u',v')\odot(u'',v'')=(u'u'',t_2)$. Again, by
Proposition~\ref{prop:pair-concatenation-least-constraint}, $t_2$ is
the least string satisfying $v'\preceq u''t_2$ and $v''\preceq t_2$.
By Lemma~\ref{lem:pair-concatenation-definedness},
$(u,v)\odot(u'u'',t_2)$ is defined iff there exists a string $t$ such
that $v\preceq u'u''t$ and $t_2\preceq t$.

If such a $t$ exists, then
$v'\preceq u''t_2\preceq u''t$ and $v''\preceq t_2\preceq t$,
so $t\in\mathcal{T}$. Conversely, if $t\in\mathcal{T}$, then
$v'\preceq u''t$ and $v''\preceq t$. By the leastness of $t_2$,
we obtain $t_2\preceq t$. Hence the right-associated concatenation is
also defined iff $\mathcal{T}\neq\emptyset$.

Its final residual is likewise the least element of $\mathcal{T}$.
Therefore, whenever either parenthesization is defined, both are
defined and both yield the same main string $uu'u''$ and the same
least residual constraint $t$. Therefore, both parenthesizations are simultaneously defined and produce the same pair. Thus, the partial concatenation operator \(\odot\) is associative.
\end{proof}

A \emph{language with string constraints} over $\Sigma$ is a subset of
$\Sigma^*\times\Sigma^*$. We write $ \mathcal{C}_{\Sigma} = \mathcal{P}(\Sigma^*\times\Sigma^*)$. For the class of all languages with string constraints over $\Sigma$.
For $R,S\in\mathcal{C}_{\Sigma}$, we define union and concatenation by
\begin{equation}
R+S = R\cup S,
\qquad
R\odot S
=
\left\{
p\odot q
\;\middle|\;
p\in R,\;
q\in S,\;
p\odot q \text{ is defined}
\right\}.
\label{eq:language-operations}
\end{equation}

The zero and identity elements are defined by $ \mathbf{0}=\emptyset, \: \mathbf{1}=\{(\varepsilon,\varepsilon)\}. $ For $R\in\mathcal{C}_{\Sigma}$, we define its positive lookahead by $
\mathsf{LA}(R)
=
\left\{
(\varepsilon,\pi(p))
\;\middle|\;
p\in R
\right\}.
$ We define the powers and Kleene star of $R$ by
$
R^0=\mathbf{1},
\:
R^{n+1}=R^n\odot R,
\:
R^*=\bigcup_{n\geq 0}R^n .
$


\begin{proposition}[\textbf{Idempotent semiring}]
\label{prop:string-constraint-idempotent-semiring}
The structure
\begin{equation}
\bigl(
\mathcal{C}_{\Sigma},
\cup,
\odot,
\mathbf{0},
\mathbf{1}
\bigr)
\label{eq:string-constraint-idempotent-semiring}
\end{equation}
is an idempotent semiring.
\end{proposition}
\begin{proof}
Set union is associative, commutative, and idempotent, with identity
$\mathbf{0}=\emptyset$. The associativity of $\odot$ follows from
Proposition~\ref{prop:pair-concatenation-associative} and its pointwise
lifting to languages. Moreover,
$R\odot\mathbf{1}=R=\mathbf{1}\odot R$ for every
$R\in\mathcal{C}_{\Sigma}$, since
$(u,v)\odot(\varepsilon,\varepsilon)=(u,v)$ and
$(\varepsilon,\varepsilon)\odot(u,v)=(u,v)$.

The empty language is absorbing for concatenation,
$R\odot\mathbf{0}=\mathbf{0}=\mathbf{0}\odot R$, and $\odot$
distributes over union on both sides by its pointwise definition.
Therefore,
$\bigl(\mathcal{C}_{\Sigma},\cup,\odot,\mathbf{0},\mathbf{1}\bigr)$
forms an idempotent semiring.
\end{proof}
The full language domain forms an idempotent semiring, while the regular fragment is generated from the basic constrained languages using union, concatenation, and Kleene star. The semantics of $\mathrm{REwPLA}$ is defined inductively by
$\mathcal{M}:\mathrm{REwPLA}\to\mathcal{C}_{\Sigma}$ as follows:
\begin{equation}
\begin{aligned}
\mathcal{M}(\emptyset)
&=\emptyset,
&
\mathcal{M}(\varepsilon)
&=\{(\varepsilon,\varepsilon)\},
&
\mathcal{M}(a)
&=\{(a,\varepsilon)\},
\\
\mathcal{M}(r+s)
&=\mathcal{M}(r)\cup\mathcal{M}(s),
&
\mathcal{M}(rs)
&=\mathcal{M}(r)\odot\mathcal{M}(s),
\\
\mathcal{M}(r^*)
&=\mathcal{M}(r)^*,
&
\mathcal{M}(\mathsf{LA}(r))
&=
\left\{
(\varepsilon,\pi(p))
\;\middle|\;
p\in\mathcal{M}(r)
\right\}.
\end{aligned}
\label{eq:rewpla-semantics}
\end{equation}
